%======================== 6. ENGINEERING ISSUES ================================
\isec{\faTools}{Engineering Issues and Lessons}

\isub{\faWifi}{Wi-Fi link: from direct AP to router mode}
\label{sec:wifi}
\begin{itemize}
  \item \textbf{Symptom:} in direct AP mode the laptop kept losing the robot (Table~\ref{tab:wifi}). The driver logs of the three AP sessions
        contain 23 \emph{Disconnected} events; in the longest session (27\,min) the link stayed up for a median of
        64.6\,s between drops (16 drops, 25 to 227\,s).
  \item \textbf{Cause on the host side:} the laptop's Realtek RTL8822CE card (driver \code{rtw_8822ce}) deauthenticated
        ``by local choice (Reason: 3=DEAUTH\_LEAVING)'', twelve times between 00:19:14 and 00:33:50, about 65\,s apart.
        Setting \code{disable_aspm=1} and \code{disable_lps_deep=1} for \code{rtw88} did not help: the next drop came
        60\,s after reconnecting (\code{notes/lab_log.md}).
  \item \textbf{Fix:} the robot was switched to router (STA) mode on the home 2.4\,GHz network, configured with a QR
        code generated by \code{tools/sta_qrcode.py} (the password is typed, never printed or saved). A 3\,min check
        (00:51:10 to 00:54:10) showed 0 deauth events, 0 driver disconnects and 0 detector re-subscribes with 10 to 15
        T3 detections/s (\code{notes/lab_log.md}); the driver log of the router-mode session then shows a single connection for 69\,min,
        covering all floor runs, until the session was ended after the final run.
  \item \textbf{Robustness kept in the code:} the detector re-subscribes to the camera after 2\,s without images, and
        the planner brakes if the tag is lost for more than 2\,s far from the goal.
\end{itemize}

\begin{table}[H]
\centering\small
\rowcolors{2}{kfupmlight}{white}
\begin{tabularx}{\linewidth}{@{}l c c c >{\raggedright\arraybackslash}X@{}}
\rowcolor{kfupm}\hd{Driver log} & \hd{Mode} & \hd{Span (min)} & \hd{Drops} & \hd{Connected time before a drop}\\
\code{driver_20261025_235452} & AP & 27.0 & 16 & median 64.6\,s, 25.0 to 227.2\,s\\
\code{driver_20261026_002315} & AP & 8.0 & 3 & median 130.1\,s, 64.1 to 263.2\,s\\
\code{driver_20261026_003846} & AP & 3.7 & 4 & median 41.1\,s, 14.0 to 64.1\,s\\
\code{driver_20261026_005040} & STA & 69.3 & 1 & one connection of 69\,min; the only drop came after the final run\\
\end{tabularx}
\caption{Connection events in the raw driver logs (\code{data/logs/}, parsed by \code{analysis/net_logs.py}).}
\label{tab:wifi}
\end{table}

\isub{\faWrench}{Other issues found on the robot}
Table~\ref{tab:issues} collects the remaining problems met during the assessment.
\begin{table}[H]
\centering\small
\rowcolors{2}{kfupmlight}{white}
\begin{tabularx}{\linewidth}{@{}>{\raggedright\arraybackslash}p{3.3cm}>{\raggedright\arraybackslash}X>{\raggedright\arraybackslash}X@{}}
\rowcolor{kfupm}\hd{Issue} & \hd{Evidence} & \hd{Resolution}\\
Arm pose hits the ToF sensor & arm hunts down to $z=0.050$\,m against the sensor (Section~\ref{sec:arm}, lab log) &
  arm set by hand 2\,cm above the sensor; planner arm guard ($\pm6$\,mm)\\
Driver $+x$ is away from the camera & 2\,s of $+0.05$\,m/s with wheels in the air: wheels turned away from the tag (lab log) & \code{drive_sign} $=-1$: \code{cmd_vel.linear.x} $=-v$; all
  logs stay in the mission frame\\
Arm TF wrong for this arm & TF gave $18^\circ$ and $22^\circ$ up; the floor tag gave $15.9^\circ$ down (lab log) & camera height and pitch from the
  floor tag (\code{extrinsics: tag})\\
Slant range at the standoff & 0.479\,m slant vs 0.362\,m ground (Figure~\ref{fig:geom}) & ground range in a level frame,
  offset within 6\,mm at 0.6 to 1.5\,m\\
Tag range limit & detected at 1.50\,m (3.8/s), ``NOT DETECTED in 8\,s'' at 1.80\,m (lab log) & runs start at 1.50\,m, the lower end of the specified
  range\\
Detection gaps at range & 4 restarts in run 1, 1.6\,s gap in run 2, 2.0\,s gap in the final run & chase tolerates gaps up
  to 2.0\,s; beyond that it searches again\\
Camera blind zone & no detection below 0.39\,m in any run & final approach on plan time and odometry\\
PyAV not installed & driver warning ``FFmpeg not available'' & raw images are enough; no change\\
Detector exit error & \code{RCLError} from a timer during shutdown, after the final run & harmless for the run; the
  planner, which owns \code{cmd_vel}, had already stopped\\
\end{tabularx}
\caption{Issues, the evidence for each, and how they were handled.}
\label{tab:issues}
\end{table}

\begin{keybox}{\faLightbulb\ \ Lessons learned}
\small
\begin{itemize}
  \item \textbf{Measure the frames on the robot itself.} Two conventions of the driver did not match this robot (drive
        direction and arm TF). A wheels-in-the-air test and one floor tag measured both directly.
  \item \textbf{A downward camera needs ground geometry.} Slant range overshoots the standoff by about 12\,cm here; the
        blind zone ends the visual feedback 0.19\,m before the stop.
  \item \textbf{Flatness makes the stop predictable.} Because the plan ends at $\dot x(T)=0$ exactly at $d_0-0.20$, the
        robot can finish the approach without the camera and still stop at 20\,cm.
  \item \textbf{Timeouts must match the sensor rate.} A 0.45\,s visibility rule suits searching, but a 4\,Hz detection
        stream at 1.5\,m needs a longer tolerance during the chase.
  \item \textbf{Fix the link before tuning the robot.} Most early failures were link drops, not control problems.
\end{itemize}
\end{keybox}
